\documentclass[12pt]{article}
\usepackage{amsmath}
\usepackage{amssymb}
\usepackage{physics}
\usepackage{hyperref}

% Shorthand for displaystyle fractions inside inline math
\newcommand{\ds}{\displaystyle}
\newcommand{\ip}[2]{\langle #1 \mid #2 \rangle}

\title{Composition of Velocities}
\author{}
\date{}

\begin{document}
\maketitle

% ============================================================
\section{Derivation of Composition of Velocities}
% ============================================================

In classical mechanics, a velocity $\vb{u}'$ in the moving frame will be measured in the
stationary frame as $\vb{u} = \vb{v} + \vb{u}'$, which is the well-known law of the
parallelograms. In relativity, the parallelogram formula holds only for velocities well below
the speed of light, and the composition of velocities needs to account for the fact that no
speed can be measured which is faster than the speed of light.

We will use the same reference systems as defined in the section on the Transformation of
Time and Space. We consider the general motion of a body in $F'$ described by the four-vector
\begin{equation}
  \ket{\vb{x}', 0} + ct'\ket{\vb{0}, 1}
\end{equation}
where we assume that $\vb{x}' = \vb{x}'(t')$. The motion of the body observed in $F$ is
given by the Lorentz transformation
\begin{equation}
  \boldsymbol{\Lambda}(\vb{v})\left(\ket{\vb{x}, 0} + ct\ket{\vb{0}, 1}\right)
  = \ket{\vb{x}', 0} + ct'\ket{\vb{0}, 1}
\end{equation}
where $\boldsymbol{\Lambda}(\vb{v})$ is the Lorentz transformation in velocity
(see \textit{The Transformation of Time and Space}, eq.~for the Lorentz transformation in
velocity):
\begin{equation}
  \boldsymbol{\Lambda}(\vb{v}) =
    \frac{\beta}{v^2}\ket{\vb{v},0}\bra{\vb{v},0}
    - \frac{\beta}{c}\ket{\vb{v},0}\bra{\vb{0},1}
    + \ket{\vb{e}_2,0}\bra{\vb{e}_2,0}
    + \ket{\vb{e}_3,0}\bra{\vb{e}_3,0}
    - \frac{\beta}{c}\ket{\vb{0},1}\bra{\vb{v},0}
    + \beta\ket{\vb{0},1}\bra{\vb{0},1}
  \label{eq:composition_Lorentz_in_velocity}
\end{equation}
and $\vb{x} = \vb{x}(t)$. The velocity of the body in $F$ is given by the derivative of the
position vector with respect to time:
\begin{equation}
  \boldsymbol{\Lambda}(\vb{v})\,\frac{d}{dt}\!\left(\ket{\vb{x},0} + ct\ket{\vb{0},1}\right)
  = \frac{d}{dt}\!\left(\ket{\vb{x}',0} + ct'\ket{\vb{0},1}\right)
\end{equation}

From the transformation of time and space (see \textit{The Transformation of Time and Space}):
\begin{equation}
  ct' = \beta\!\left(-\frac{v}{c}\ip{\vb{e}_1}{\vb{x}} + ct\right)
\end{equation}
we can compute the relation between the velocities $\vb{u}$ and $\vb{u}'$ in $F$ and $F'$
respectively,
\begin{equation}
  \boldsymbol{\Lambda}(\vb{v})\bigl(\ket{\vb{u},0} + c\ket{\vb{0},1}\bigr)
  = \beta\!\left(1 - \frac{\ip{\vb{v}}{\vb{u}}}{c^2}\right)
    \bigl(\ket{\vb{u}',0} + c\ket{\vb{0},1}\bigr)
\end{equation}
which, using the Lorentz transformation in velocity \eqref{eq:composition_Lorentz_in_velocity}, reads
\begin{equation}
  \begin{aligned}
    &\beta\!\left(\frac{\ip{\vb{v}}{\vb{u}}}{v^2} - 1\right)\ket{\vb{v},0}
    + \ket{\vb{e}_2,0}\ip{\vb{e}_2}{\vb{u}}
    + \ket{\vb{e}_3,0}\ip{\vb{e}_3}{\vb{u}}
    + \ket{\vb{0},1}\,\beta c\!\left(1 - \frac{\ip{\vb{v}}{\vb{u}}}{c^2}\right) \\
    &= \beta\!\left(1 - \frac{\ip{\vb{v}}{\vb{u}}}{c^2}\right)
       \bigl(\ket{\vb{u}',0} + c\ket{\vb{0},1}\bigr)
  \end{aligned}
\end{equation}
and simplifies to,
\begin{equation}
  \beta\!\left(\frac{\ip{\vb{v}}{\vb{u}}}{v^2} - 1\right)\ket{\vb{v},0}
  + \ket{\vb{e}_2,0}\ip{\vb{e}_2}{\vb{u}}
  + \ket{\vb{e}_3,0}\ip{\vb{e}_3}{\vb{u}}
  = \beta\!\left(1 - \frac{\ip{\vb{v}}{\vb{u}}}{c^2}\right)\ket{\vb{u}',0}
\end{equation}

The equation above can be written in a more compact form as,
\begin{equation}
  \beta\,\frac{\ip{\vb{v}}{\vb{u}}}{v^2}\ket{\vb{v},0}
  - \beta\ket{\vb{v},0}
  - \frac{\ip{\vb{v}}{\vb{u}}}{v^2}\ket{\vb{v},0}
  + \ket{\vb{u}}
  = \beta\!\left(1 - \frac{\ip{\vb{v}}{\vb{u}}}{c^2}\right)\ket{\vb{u}',0}
\end{equation}
or,
\begin{equation}
  (\beta - 1)\frac{\ip{\vb{v}}{\vb{u}}}{v^2}\ket{\vb{v},0}
  - \beta\ket{\vb{v},0}
  + \ket{\vb{u}}
  = \beta\!\left(1 - \frac{\ip{\vb{v}}{\vb{u}}}{c^2}\right)\ket{\vb{u}',0}
  \label{eq:velocities_relationship}
\end{equation}
from which we readily obtain the velocity $\ket{\vb{u}'}$ as,
\begin{equation}
  \ket{\vb{u}',0} = \frac{1}{\displaystyle 1 - \frac{\ip{\vb{v}}{\vb{u}}}{c^2}}
  \left(
    \frac{\beta-1}{\beta}\,\frac{\ip{\vb{v}}{\vb{u}}}{v^2}\ket{\vb{v},0}
    - \ket{\vb{v},0}
    + \frac{\ket{\vb{u},0}}{\beta}
  \right).
  \label{eq:inverse_velocity_composition}
\end{equation}

The velocity $\ket{\vb{u}}$ can be easily expressed in terms of the velocity $\ket{\vb{u}'}$
by exchanging the roles of $\ket{\vb{u}}$ and $\ket{\vb{u}'}$ in the equation above, and
changing the sign of the velocity.

As an exercise, we will derive the same result in a more general way by inverting the equation
above. We start by factoring the velocity $\ket{\vb{u}}$ in \eqref{eq:velocities_relationship}
and express the equation as an operator on the velocity $\ket{\vb{u}}$,
\begin{equation}
  \left(
    \frac{\beta}{v^2}\ket{\vb{v}}\bra{\vb{v}}
    + \ket{\vb{e}_2,0}\bra{\vb{e}_2,0}
    + \ket{\vb{e}_3,0}\bra{\vb{e}_3,0}
    + \frac{\beta}{c^2}\ket{\vb{u}'}\bra{\vb{v}}
  \right)\ket{\vb{u}}
  = \beta\bigl(\ket{\vb{u}'} + \ket{\vb{v}}\bigr)
\end{equation}

Further expanding the equation above in the components of $\ket{\vb{u}'}$, we get
\begin{equation}
  \begin{aligned}
    \Bigl(
      &\beta\!\left(1 + \frac{v}{c^2}u'_1\right)\ket{\vb{e}_1,0}\bra{\vb{e}_1,0}
      + \beta\,\frac{v}{c^2}u'_2\,\ket{\vb{e}_2,0}\bra{\vb{e}_1,0}
      + \beta\,\frac{v}{c^2}u'_3\,\ket{\vb{e}_3,0}\bra{\vb{e}_1,0} \\
      &+ \ket{\vb{e}_2,0}\bra{\vb{e}_2,0}
      + \ket{\vb{e}_3,0}\bra{\vb{e}_3,0}
    \Bigr)\cdot\ket{\vb{u},0}
    = \beta\bigl(\ket{\vb{u}'} + \ket{\vb{v}}\bigr)
  \end{aligned}
\end{equation}

The operator,
\begin{equation}
  \beta\!\left(1 + \frac{v}{c^2}u'_1\right)\ket{\vb{e}_1,0}\bra{\vb{e}_1,0}
  + \beta\,\frac{v}{c^2}u'_2\,\ket{\vb{e}_2,0}\bra{\vb{e}_1,0}
  + \beta\,\frac{v}{c^2}u'_3\,\ket{\vb{e}_3,0}\bra{\vb{e}_1,0}
  + \ket{\vb{e}_2,0}\bra{\vb{e}_2,0}
  + \ket{\vb{e}_3,0}\bra{\vb{e}_3,0}
\end{equation}
can be easily inverted as,
\begin{equation}
  \begin{aligned}
    &\frac{1}{\displaystyle\beta\!\left(1 + \frac{v}{c^2}u'_1\right)}
      \ket{\vb{e}_1,0}\bra{\vb{e}_1,0}
    - \frac{v}{c^2}\frac{u'_2}{\displaystyle 1 + \frac{v}{c^2}u'_1}
      \ket{\vb{e}_2,0}\bra{\vb{e}_1,0} \\
    &- \frac{v}{c^2}\frac{u'_3}{\displaystyle 1 + \frac{v}{c^2}u'_1}
      \ket{\vb{e}_3,0}\bra{\vb{e}_1,0}
    + \ket{\vb{e}_2,0}\bra{\vb{e}_2,0}
    + \ket{\vb{e}_3,0}\bra{\vb{e}_3,0}
  \end{aligned}
\end{equation}

We can express the velocity $\ket{\vb{u}}$ as,
\begin{equation}
  \begin{aligned}
    \ket{\vb{u},0} =\;
    &\left(
      \frac{1}{\displaystyle\beta\!\left(1 + \frac{v}{c^2}u'_1\right)}
        \ket{\vb{e}_1,0}\bra{\vb{e}_1,0}
      - \frac{v}{c^2}\frac{u'_2}{\displaystyle 1 + \frac{v}{c^2}u'_1}
        \ket{\vb{e}_2,0}\bra{\vb{e}_1,0} \right.\\
    &\left.
      - \frac{v}{c^2}\frac{u'_3}{\displaystyle 1 + \frac{v}{c^2}u'_1}
        \ket{\vb{e}_3,0}\bra{\vb{e}_1,0}
      + \ket{\vb{e}_2,0}\bra{\vb{e}_2,0}
      + \ket{\vb{e}_3,0}\bra{\vb{e}_3,0}
    \right) \\[0.4em]
    &\cdot\left(v\beta\ket{\vb{e}_1,0} + \beta\ket{\vb{u}'}\right)
  \end{aligned}
\end{equation}
which finally simplifies to,
\begin{equation}
  \ket{\vb{u},0} = \frac{1}{\displaystyle 1 + \frac{v}{c^2}u'_1}
  \left(
    (v + u'_1)\ket{\vb{e}_1,0}
    - \sqrt{1 - \frac{v^2}{c^2}}\,u'_2\,\ket{\vb{e}_2,0}
    - \sqrt{1 - \frac{v^2}{c^2}}\,u'_3\,\ket{\vb{e}_3,0}
  \right)
\end{equation}

By dropping the explicit time dependence, we can write the above equation in a more elegant
form as,
\begin{equation}
  \ket{\vb{u}} =
  \frac{1}{\displaystyle 1 + \frac{\ip{\vb{v}}{\vb{u}'}}{c^2}}
  \left(
    \frac{\beta - 1}{\beta}\,\frac{\ip{\vb{u}'}{\vb{v}}}{v^2}\ket{\vb{v}}
    + \ket{\vb{v}}
    + \frac{\ket{\vb{u}'}}{\beta}
  \right)
  \label{eq:velocity_composition}
\end{equation}

% ============================================================
\section{Limit Formulas}
% ============================================================

In this section we are going to derive the limit formulas for the composition of velocities
when one of the velocities approaches the speed of light. We start by considering the limit
of the velocity composition formula \eqref{eq:velocity_composition} when the velocity
$\ket{\vb{u}'}$ approaches the speed of light.

\subsection{$\ket{\vb{u}}$ Approaching Speed of Light}

We start by first considering the case where $\ket{\vb{u}}$ approaches the speed of light
$c$ and we will determine the value that $\ket{\vb{u}'}$ reaches. To achieve this, we need
to compute the modulus of $\ket{\vb{u}'}$ from \eqref{eq:inverse_velocity_composition},
which reads:
\begin{equation}
  u'^2 = \frac{1}{\displaystyle\left(1 - \frac{\ip{\vb{v}}{\vb{u}}}{c^2}\right)^2}
  \left(
    \frac{\ip{\vb{u}}{\vb{v}}^2}{c^2} + v^2 + \frac{u^2}{\beta^2}
    - 2\ip{\vb{u}}{\vb{v}}
  \right)
\end{equation}

When $\ket{\vb{u}} = c\ket{\vb{n}}$, where $\ket{\vb{n}}$ is a unit vector, the equation
above simplifies to
\begin{equation}
  u'^2 = \frac{c^2}{\displaystyle\left(1 - \frac{\ip{\vb{v}}{\vb{n}}}{c}\right)^2}
  \left(1 - \frac{\ip{\vb{v}}{\vb{n}}}{c}\right)^2 = c^2
\end{equation}
which implies that the velocity $\ket{\vb{u}'}$ reaches the speed of light when $\ket{\vb{u}}$
approaches the speed of light.

\subsection{$\ket{\vb{u}'}$ Approaching Speed of Light}

Similarly, we can determine the limit of the velocity $\ket{\vb{u}}$ when the velocity
$\ket{\vb{u}'}$ approaches the speed of light. From the velocity composition formula
\eqref{eq:velocity_composition}, the modulus of the velocity $\ket{\vb{u}}$ reads,
\begin{equation}
  u^2 = \frac{1}{\displaystyle\left(1 + \frac{\ip{\vb{v}}{\vb{u}'}}{c^2}\right)^2}
  \left(
    \frac{\ip{\vb{u}'}{\vb{v}}^2}{v^2} + v^2 + \frac{u'^2}{\beta^2}
    + 2\ip{\vb{u}'}{\vb{v}}
  \right)
\end{equation}
which is the same equation for the modulus of $\ket{\vb{u}}$ with the sign of the velocity
$\ket{\vb{v}}$ reversed. When $\ket{\vb{u}'} = c\ket{\vb{n}}$, where $\ket{\vb{n}}$ is a
unit vector, the equation above simplifies to
\begin{equation}
  u^2 = \frac{c^2}{\displaystyle\left(1 + \frac{\ip{\vb{v}}{\vb{n}}}{c}\right)^2}
  \left(1 + \frac{\ip{\vb{v}}{\vb{n}}}{c}\right)^2 = c^2
\end{equation}
which, also in this case, means that the velocity $\ket{\vb{u}}$ reaches the speed of light
when $\ket{\vb{u}'}$ approaches the speed of light.

\subsection{$\ket{\vb{v}}$ Approaching Speed of Light}

Following the same logic, when the velocity $\ket{\vb{v}} = c\ket{\vb{n}}$, the modulus of
the velocity $\ket{\vb{u}}$ reads,
\begin{equation}
  u^2 = \frac{c^2}{\displaystyle\left(1 + \frac{\ip{\vb{u}'}{\vb{n}}}{c}\right)^2}
  \left(1 + \frac{\ip{\vb{u}'}{\vb{n}}}{c}\right)^2 = c^2
\end{equation}
independently of the value of the velocity $\ket{\vb{u}'}$.

\end{document}